\section{Qwen3-Coder：Agentic Coding in the World}

前一章看通用产品 Agent，本章看代码环境里的 Agent，核心问题是：为什么 coding 是最适合 Agent 训练的高价值环境之一，以及开源模型如何在这个环境里追赶。第三份材料是 Qwen3-Coder。官方博文介绍 Qwen3-Coder-480B-A35B-Instruct：480B MoE、35B active，原生 256K context，通过 YaRN 可扩到 1M。它面向 agentic coding，强调 code RL、long-horizon RL、20,000 并行环境，以及 Qwen Code 工具链。

\begin{figure}[H]
\centering
\includegraphics[width=0.96\textwidth]{qwen3-coder-stack.png}
\caption{Qwen3-Coder Stack：480B MoE、256K/1M context、Code RL 与 Long-horizon RL。自制概念图，依据 Qwen 官方博文和 01:53:38--01:59:04 对谈内容整理。}
\end{figure}

\begin{knowledgebox}{读图：代码是最适合 Agent RL 的环境之一}
代码任务 hard to solve, easy to verify：解决很难，但测试和执行反馈相对明确。这让代码成为大规模 RL 和 long-horizon interaction 的好环境。
\end{knowledgebox}

\subsection{Code RL 与 Long-horizon RL}

上一节给出 Qwen3-Coder 的模型栈，本节解释它为什么强调 RL。Qwen3-Coder 博文强调，代码任务天然适合 execution-driven RL。真实软件工程任务需要多轮交互：规划、使用工具、接收反馈、修正决策。Qwen3-Coder 引入 long-horizon RL，并构建可并行运行的大规模环境。

\begin{importantbox}{为什么 coding 是 Agent 的高地}
代码环境有清晰动作、可执行反馈、测试验证和高价值任务。相比很多开放世界任务，coding 更容易定义 reward，也更容易规模化训练和评测。
\end{importantbox}

\begin{knowledgebox}{Code RL 为什么“难解但易验”}
很多代码任务的解决过程很难，需要理解代码库、规划修改和调试；但结果可以通过测试、编译、lint、benchmark 或用户脚本验证。这种结构非常适合强化学习，因为 reward 可以相对自动化。
\end{knowledgebox}

\subsection{本章小结}

Qwen3-Coder 展示了开源模型在 agentic coding 上的系统路线：长上下文、代码数据、执行驱动 RL、长程工具交互和 CLI 工具链。

